\specsection{一般拓扑空间上的测度与紧收敛}

对拓扑与测度论之间联系的研究，首先被视为对测度正则性性质、特别是“外”正则与“内”正则的研究；\(^{(14)}\) 在无穷处可数的局部紧空间上，内正则等价于外正则。Lebesgue 对直线上子集测度的构造突出了这两种正则性，而波兰空间上测度的外正则性质似乎在 1935 年前后已为公众所知。但直到 1940 年，A. D. Alexandroff 才在一篇传播因战争而推迟的文章 (XIX) 中突出了内正则的作用，并证明波兰空间上的测度具有这一性质；这一结果后来被 Prokhorov (XIII) 重新发现，并且常被错误地归功于这位作者。直到最近人们才认识到这一性质可以推广到 Souslin 空间；因此，这些空间的重要性大大增加，尤其是因为人们认识到它们的理论可以在没有任何可度量化假设的情况下展开，而且几乎全部的函数空间都是 Souslin 的（多数情形甚至是 Lusin 的）。\(^{(15)}\) 这些就是促使我们在本章强调内正则测度的理由。

测度收敛方式（淡收敛或紧收敛）的定义，最方便的做法是把测度空间与一个连续函数空间做成对偶。A. A. Markoff 推广了 F. Riesz 的一个老结果，于 1938 年建立了 \(\mathcal{C}(\mathrm{X})\) 上正泛函与紧空间 X 上正则测度之间的一一对应。在前面已引用的著作 (XIX) 中，A. D. Alexandroff 把这些结果推广到完全正则空间的情形：他在空间 \(\mathcal{C}^b (\mathrm{X})\)（完全正则空间 \(\mathrm{X},^{(16)}\) 上有界连续函数的）上的正线性形式所成的集合中引入了一个层级，他定义了有界测度的紧收敛，并证明了下面两个定理：

\begin{enumerate}
\def\labelenumi{\alph{enumi})}
\item
  若 X 是波兰空间，则对应于测度的 \(\mathcal{C}^{b}(X)\) 上线性形式所成的集合对序列的弱收敛是闭的；
\item
  若一列有界测度有紧极限，则“没有质量逃逸到无穷远”（这是关于紧收敛的 Prokhorov 定理之逆的一个弱形式）。
\end{enumerate}

从这些丰富的概念与定理中，Prokhorov 提取了对随机过程理论重要的结果，并以简单而醒目的形式呈现它们。在他 1956 年已引用的大著作 (XIII) 中，很大一部分专门讨论波兰空间上的有界正测度；他推广了 Lévy 的一个构造，在质量为 1 的正测度所成的集合上定义了一个距离，使它成为一个波兰空间，然后建立了一个重要的紧性判别法（cf.~§5, No.~5, Th. 1）。与 Prokhorov 相独立，Le Cam (XX) 得到了若干关于测度紧收敛的紧性结果；他对所考虑的空间不作任何可度量化假设，而且在局部紧情形他的结果归结为 Dieudonné 的较早的定理。

